% Бүлэг 4

\chapter{Хөгжүүлэлт ба туршилт}
\label{Chapter4}
\pagecolor{white}

\section{Системийн архитектур}

Системийн ерөнхий архитектурыг гурван давхаргат загварын дагуу зохион байгуулсан. Backend-ийг Routes → Controller → Service → Repository → Worker давхрагуудтайгаар хөгжүүлсэн.

\begin{table}[h!]
\centering
\caption{Хэрэгжүүлсэн технологийн стек}
\label{tab:tech_stack}
\begin{tabular}{|l|l|p{5cm}|}
\hline
\textbf{Давхарга} & \textbf{Технологи} & \textbf{Зориулалт} \\
\hline
Frontend & React 18, TypeScript, Vite & SPA интерфэйс \\
\hline
Backend & Node.js, Express 5, TypeScript & RESTful API сервер \\
\hline
Өгөгдлийн сан & PostgreSQL, Sequelize ORM & Бүтцэт өгөгдөл хадгалах \\
\hline
Медиа хадгалалт & Cloudinary CDN & Видео, зураг, аудио файл \\
\hline
Скрипт AI & Anthropic Claude, Groq/Llama 3 & Сценарий бичих \\
\hline
Зураг AI & Google Gemini Flash Image 2.5 & Scene зураг үүсгэх \\
\hline
Аудио AI & ElevenLabs, Chimege, Gemini TTS & Хоолой дуурайлгах \\
\hline
Видео угсралт & FFmpeg & MP4 нэгтгэлт \\
\hline
Бодит цаг & SSE (Server-Sent Events) & Явц дамжуулах \\
\hline
Баталгаажуулалт & JWT, bcrypt & Хэрэглэгч нэвтрэх \\
\hline
\end{tabular}
\end{table}

\section{Өгөгдлийн сангийн зохиомж}

\subsection{Өгөгдлийн сангийн схемийн ерөнхий бүтэц}

Өгөгдлийн сангийн схемийг 8 үндсэн хүснэгт, 7 төрлийн холбоос, 8 индекстэйгээр зохиомжилж 3NF хүртэл нормчилсан. Үндсэн хүснэгтүүд: users, videos, scenes, subscriptions, projects, admins.

\subsection{users хүснэгт}

\code{users} хүснэгт нь системийн бүртгэлтэй хэрэглэгчдийн мэдээллийг хадгална.

\begin{table}[h!]
\centering
\caption{users хүснэгтийн схем}
\label{tab:users}
\begin{tabular}{|l|l|l|p{4cm}|}
\hline
\textbf{Талбар} & \textbf{Төрөл} & \textbf{Хязгаарлалт} & \textbf{Тайлбар} \\
\hline
id & UUID & PRIMARY KEY & Хэрэглэгчийн ID \\
\hline
email & VARCHAR(255) & UNIQUE, NOT NULL & Нэвтрэх имэйл \\
\hline
password & VARCHAR(255) & NOT NULL & Bcrypt хэш \\
\hline
name & VARCHAR(100) & NOT NULL & Хэрэглэгчийн нэр \\
\hline
role & ENUM & DEFAULT 'user' & user / pro / manager \\
\hline
is\_active & BOOLEAN & DEFAULT true & Идэвхтэй эсэх \\
\hline
created\_at & TIMESTAMP & DEFAULT now() & Бүртгэлийн огноо \\
\hline
updated\_at & TIMESTAMP & DEFAULT now() & Шинэчилсэн огноо \\
\hline
\end{tabular}
\end{table}

\subsection{videos хүснэгт}

\code{videos} хүснэгт нь үүсгэгдсэн бүх видеоны метадата, статус, үр дүнгийн холбоосыг хадгална.

\begin{table}[h!]
\centering
\caption{videos хүснэгтийн схем}
\label{tab:videos}
\begin{tabular}{|l|l|l|p{4cm}|}
\hline
\textbf{Талбар} & \textbf{Төрөл} & \textbf{Хязгаарлалт} & \textbf{Тайлбар} \\
\hline
id & UUID & PRIMARY KEY & Видео ID \\
\hline
user\_id & UUID & FOREIGN KEY & Эзэмшигч хэрэглэгч \\
\hline
project\_id & UUID & FOREIGN KEY, NULL & Харьяалагдах төсөл \\
\hline
title & VARCHAR(200) & NOT NULL & Видеоны гарчиг \\
\hline
topic & TEXT & NOT NULL & Видеоны сэдэв \\
\hline
genre & VARCHAR(50) & NOT NULL & Видеоны жанр \\
\hline
language & VARCHAR(20) & NOT NULL & Хэл (mn / en) \\
\hline
status & ENUM & DEFAULT 'pending' & pending / processing / done / failed \\
\hline
video\_url & TEXT & NULL & Cloudinary-д байршсан URL \\
\hline
duration & FLOAT & NULL & Видеоны урт (секунд) \\
\hline
file\_size & BIGINT & NULL & Видео файлын хэмжээ \\
\hline
created\_at & TIMESTAMP & DEFAULT now() & Үүсгэсэн огноо \\
\hline
\end{tabular}
\end{table}

\subsection{scenes хүснэгт}

\code{scenes} хүснэгт нь видеоны хэсэг бүрийн narration текст, image prompt, дараалал зэргийг хадгална. Studio засварлагч энэ хүснэгтэн дээр голчлон ажилладаг.

\begin{table}[h!]
\centering
\caption{scenes хүснэгтийн схем}
\label{tab:scenes}
\begin{tabular}{|l|l|l|p{4cm}|}
\hline
\textbf{Талбар} & \textbf{Төрөл} & \textbf{Хязгаарлалт} & \textbf{Тайлбар} \\
\hline
id & UUID & PRIMARY KEY & Scene ID \\
\hline
video\_id & UUID & FOREIGN KEY & Харьяалагдах видео \\
\hline
scene\_index & INTEGER & NOT NULL & Дараалал (0-аас) \\
\hline
narration & TEXT & NOT NULL & Хэлэх текст \\
\hline
image\_prompt & TEXT & NOT NULL & Зураг үүсгэх prompt \\
\hline
image\_url & TEXT & NULL & Үүсгэгдсэн зураг URL \\
\hline
audio\_url & TEXT & NULL & TTS аудио URL \\
\hline
audio\_duration & FLOAT & NULL & Аудионы үргэлжлэх хугацаа \\
\hline
tts\_provider & VARCHAR(50) & NOT NULL & gemini / elevenlabs / chimege \\
\hline
voiceId & VARCHAR(100) & NOT NULL & Дуу хоолойн ID \\
\hline
transition\_type & VARCHAR(50) & NULL & Scene хоорондын шилжилт \\
\hline
motion\_effect & VARCHAR(50) & NULL & Зурагны хөдөлгөөний эффект \\
\hline
created\_at & TIMESTAMP & DEFAULT now() & Үүсгэсэн огноо \\
\hline
\end{tabular}
\end{table}

\subsection{Хүснэгтүүдийн хоорондын холбоос}

\begin{table}[h!]
\centering
\caption{Хүснэгтүүдийн хоорондын холбоосын матриц}
\label{tab:relations}
\begin{tabular}{|l|l|c|p{4cm}|}
\hline
\textbf{Эх хүснэгт} & \textbf{Зорилтот хүснэгт} & \textbf{Төрөл} & \textbf{Тайлбар} \\
\hline
users & videos & 1:N & Нэг хэрэглэгч олон видеотой \\
\hline
users & projects & 1:N & Нэг хэрэглэгч олон төсөлтэй \\
\hline
projects & videos & 1:N & Нэг төсөл олон видеотой \\
\hline
videos & scenes & 1:N & Нэг видео олон scene-тэй \\
\hline
users & subscriptions & 1:1 & Нэг хэрэглэгч нэг subscription \\
\hline
\end{tabular}
\end{table}

\subsection{Индекс ба гүйцэтгэлийн оптимизаци}

Хайлт болон жагсаалтын хурдыг сайжруулахын тулд дараах индексүүдийг үүсгэв:

\begin{itemize}
\item \code{users.email} --- UNIQUE INDEX (нэвтрэхэд хайх)
\item \code{videos.user\_id} --- INDEX (хэрэглэгчийн видеог жагсаахад)
\item \code{videos.project\_id} --- INDEX (төслийн видеог жагсаахад)
\item \code{videos.status} --- INDEX (processing төлөвтэй видеог шүүхэд)
\item \code{videos.created\_at} --- INDEX (огноогоор эрэмбэлэхэд)
\item \code{scenes.video\_id} --- INDEX (нэг видеоны scene-уудыг авахад)
\item \code{scenes.video\_id + scene\_index} --- COMPOSITE INDEX (дараалсан хайлтад)
\end{itemize}

\section{Backend API хөгжүүлэх}

Backend API-ийг Node.js v22, Express v4, TypeScript ашиглан хөгжүүлсэн. RESTful архитектурын зарчмын дагуу resource-based endpoint бүтцийг ашигласан.

\subsection{Технологийн стек}

\begin{table}[h!]
\centering
\caption{Backend-д ашигласан технологийн стек}
\label{tab:backend_stack}
\begin{tabular}{|l|p{7cm}|}
\hline
\textbf{Технологи / Сан} & \textbf{Зориулалт} \\
\hline
Node.js v22 & JavaScript runtime \\
\hline
Express v4 & Web framework, routing \\
\hline
TypeScript v5 & Static type checking \\
\hline
Sequelize v6 & PostgreSQL ORM \\
\hline
jsonwebtoken & JWT баталгаажуулалт \\
\hline
bcrypt & Нууц үг хэшлэх \\
\hline
zod & Request payload validation \\
\hline
winston & Structured logging \\
\hline
express-rate-limit & Rate limiting middleware \\
\hline
helmet & HTTP security headers \\
\hline
multer & File upload боловсруулалт \\
\hline
socket.io & Real-time progress updates \\
\hline
fluent-ffmpeg & FFmpeg-ийн Node.js binding \\
\hline
\end{tabular}
\end{table}

\subsection{API endpoint жагсаалт}

Системийн backend нь дараах REST endpoint-уудаас бүрдэнэ. Бүх endpoint \code{/api/v1} префикстэй.

\begin{table}[h!]
\centering
\caption{Authentication endpoint-ууд}
\label{tab:auth_api}
\begin{tabular}{|l|l|p{5cm}|}
\hline
\textbf{Method} & \textbf{Endpoint} & \textbf{Тайлбар} \\
\hline
POST & /auth/register & Шинэ хэрэглэгч бүртгэх \\
\hline
POST & /auth/login & Нэвтрэх ба JWT token авах \\
\hline
POST & /auth/refresh & Refresh token-оос шинэ access token авах \\
\hline
POST & /auth/logout & Гарах ба токен цуцлах \\
\hline
POST & /auth/forgot-password & Нууц үг сэргээх имэйл илгээх \\
\hline
POST & /auth/reset-password & Нууц үг шинэчлэх \\
\hline
GET & /auth/me & Одоогийн хэрэглэгчийн мэдээлэл авах \\
\hline
\end{tabular}
\end{table}

\begin{table}[h!]
\centering
\caption{Video endpoint-ууд}
\label{tab:video_api}
\begin{tabular}{|l|l|p{5cm}|l|}
\hline
\textbf{Method} & \textbf{Endpoint} & \textbf{Тайлбар} & \textbf{Хамгаалалт} \\
\hline
POST & /api/video/generate & Видео үүсгэх & JWT \\
\hline
GET & /api/video/progress/:jobId & SSE явц & --- \\
\hline
GET & /api/video & Хэрэглэгчийн видеонууд & JWT \\
\hline
GET & /api/video/:videoId & Видеоны дэлгэрэнгүй & JWT \\
\hline
DELETE & /api/video/:videoId & Видео устгах & JWT \\
\hline
POST & /api/video/:videoId/regenerate/:sceneIndex & Scene медиа дахин үүсгэх & JWT \\
\hline
POST & /api/video/:videoId/reassemble & Видеог дахин угсрах & JWT \\
\hline
POST & /api/video/:videoId/regen-text & Scene текст AI-аар дахин бичих & JWT \\
\hline
\end{tabular}
\end{table}

\begin{table}[h!]
\centering
\caption{Аудио endpoint-ууд}
\label{tab:audio_api}
\begin{tabular}{|l|l|p{6cm}|}
\hline
\textbf{Method} & \textbf{Endpoint} & \textbf{Тайлбар} \\
\hline
POST & /api/audio/generate & TTS аудио үүсгэх (ElevenLabs/Chimege) \\
\hline
POST & /api/audio/generate-gemini & Gemini TTS аудио үүсгэх \\
\hline
POST & /api/audio/generate-from-script & Скриптийн бүх scene-д аудио үүсгэх \\
\hline
\end{tabular}
\end{table}

\subsection{Authentication ба JWT хэрэгжүүлэлт}

Баталгаажуулалтыг JSON Web Token (JWT) дээр суурилуулсан. Хэрэглэгч нэвтрэх бүрд хоёр төрлийн токен олгоно:

\begin{enumerate}
\item \textbf{Access token} --- 15 минутын хугацаатай, бүх API хүсэлтэд Authorization header-аар дамжуулна.
\item \textbf{Refresh token} --- 7 хоногийн хугацаатай, HttpOnly cookie-д хадгалагдаж, access token дуусах үед шинэ access token авахад ашиглагдана.
\end{enumerate}

Энэ хослол нь хоёр давуу талтай. Access token хулгайлагдсан тохиолдолд 15 минутаас илүү ашиглах боломжгүй. Refresh token нь HttpOnly cookie-д хадгалагдсан тул JavaScript-ээс хандаж болохгүй учир XSS халдлагаас хамгаалагдана.

\subsection{Middleware ба эрхийн шалгалт}

Систем нь дараах middleware давхаргуудтай:

\begin{itemize}
\item \textbf{authMiddleware} --- JWT token шалгах, хэрэглэгчийн мэдээллийг request-д хавсаргах
\item \textbf{roleMiddleware} --- хэрэглэгчийн role-г шалгаж зөвшөөрөл олгох/татгалзах
\item \textbf{rateLimitMiddleware} --- хэт олон хүсэлтийг хязгаарлах
\item \textbf{validationMiddleware} --- Zod схемээр request payload шалгах
\item \textbf{errorHandlerMiddleware} --- алдааг нэгдсэн байдлаар боловсруулж JSON хариу буцаах
\end{itemize}

\section{Frontend интерфэйс хөгжүүлэх}

\subsection{Технологийн стек}

\begin{table}[h!]
\centering
\caption{Frontend технологийн стек}
\label{tab:frontend_stack}
\begin{tabular}{|l|p{7cm}|}
\hline
\textbf{Технологи} & \textbf{Зориулалт} \\
\hline
React 18 & UI component framework \\
\hline
TypeScript & Static type safety \\
\hline
Vite & Build tool, dev server \\
\hline
Tailwind CSS & Utility-first styling \\
\hline
React Router v6 & Client-side routing \\
\hline
Zustand & Global state management \\
\hline
Axios & HTTP client \\
\hline
React Query & Server state, caching \\
\hline
Framer Motion & Animation \\
\hline
\end{tabular}
\end{table}

\subsection{Component ба хуудсын архитектур}

Системийн frontend нь дараах үндсэн хуудсуудаас бүрдэнэ:

\begin{itemize}
\item \textbf{Landing Page} --- танилцуулга, нэвтрэх, бүртгүүлэх
\item \textbf{Dashboard} --- хэрэглэгчийн нүүр хуудас, видеоны жагсаалт
\item \textbf{Generate Page} --- видео үүсгэх форм болон явцын харуулагч
\item \textbf{Studio Page} --- AI Video Studio засварлагч
\item \textbf{Projects Page} --- folder зохион байгуулалт
\item \textbf{Admin Dashboard} --- хэрэглэгч, видео, статистик удирдлага
\end{itemize}

\subsection{AI Video Studio засварлагчийн UI}

Studio засварлагч нь дараах 4 үндсэн хэсгээс бүрдэнэ:

\begin{enumerate}
\item \textbf{Scene timeline} --- видеоны бүх scene-г дарааллан харуулах, сонгох
\item \textbf{Preview player} --- сонгосон scene-г preview болон эцсийн видеог үзэх
\item \textbf{Editor panel} --- narration текст засварлах, image prompt өөрчлөх, дуу хоолой сонгох
\item \textbf{Re-render status bar} --- дахин үүсгэлтийн явц хянах
\end{enumerate}

\section{AI service-үүдтэй холболт}

\subsection{Groq + Llama 3.3 70B (скрипт үүсгэлт)}

Groq API-г ашиглан хэрэглэгчийн оруулсан сэдвээс видеоны скриптийг үүсгэнэ. Скрипт нь scene бүрийн дараах мэдээллийг агуулна: scene дугаар, цагийн тэмдэглэл, богино тайлбар, narration текст, зураг үүсгэх image prompt.

Groq-г сонгосон үндэслэл нь LPU чипс дээр ажиллах тул маш хурдан (токен/секунд өндөр), Llama 3.3 70B нь бүтэцтэй JSON гаралт сайн дагадаг, өртөг харьцангуй бага байдаг.

\subsection{Imagen 4 / Nano Banana (зураг үүсгэлт)}

Google Gemini Flash Image 2.5-г scene бүрийн зураг үүсгэхэд ашиглана. Image prompt нь scene-ийн narration-аас автоматаар гарна. Хурд болон чанарын тэнцвэрийн хувьд Gemini Flash Image 2.5 нь production орчинд тохиромжтой.

\subsection{ElevenLabs (дуу хоолой синтез)}

ElevenLabs Flash v2.5 болон Multilingual v2 загваруудыг scene бүрийн narration-г дуу болгон хөрвүүлэхэд ашиглана. ElevenLabs-г сонгосон шалтгаан: байгалийн дуу хоолой, олон хэл дэмжлэг, харьцангуй бага latency.

Монгол хэлний fallback-д Chimege API болон Google Gemini TTS-г ашиглана.

\subsection{FFmpeg видео угсралтын pipeline}

FFmpeg нь видео угсралтын үндсэн хэрэгсэл болно. Pipeline нь дараах алхмуудтай:

\begin{enumerate}
\item Scene бүрийн зураг + аудио + хугацааг нийлэгжүүлэх
\item Зөвлөмж бүрт хөдөлгөөний эффект (Ken Burns) нэмэх
\item Scene хоорондын шилжилт хийх
\item Subtitle (SRT) нэмэх
\item Background music нэмэх
\item Эцсийн MP4 файл гаргах
\end{enumerate}

\subsection{Cloudinary болон AWS S3 нийлмэл стратеги}

Бэлэн видео болон медиа файлуудыг Cloudinary CDN-д байршуулна. Том файл болон нөөц хадгалалтад AWS S3-г ашиглана. Cloudinary нь автоматар format оптимизаци, adaptive streaming боломжтой.

\section{Системийн туршилт}

\subsection{Туршилтын стратеги}

Туршилтын хувьд unit, integration, API, E2E, гүйцэтгэлийн туршилт зэрэг олон давхаргат стратегийг хэрэгжүүллэв:

\begin{table}[h!]
\centering
\caption{Туршилтын стратегийн пирамид}
\label{tab:test_strategy}
\begin{tabular}{|l|l|c|c|}
\hline
\textbf{Туршилтын төрөл} & \textbf{Хэрэгсэл} & \textbf{Тооны хувь} & \textbf{Хурд} \\
\hline
Unit tests & Jest, Vitest & 60\% & Маш хурдан \\
\hline
Integration tests & Jest + Supertest & 25\% & Хурдан \\
\hline
API tests & Postman, Newman & 10\% & Дунд \\
\hline
E2E tests & Playwright & 5\% & Удаан \\
\hline
\end{tabular}
\end{table}

\subsection{Туршилтын хэрэгсэл}

\begin{table}[h!]
\centering
\caption{Туршилтын хэрэгсэл}
\label{tab:test_tools}
\begin{tabular}{|l|p{8cm}|}
\hline
\textbf{Хэрэгсэл} & \textbf{Зориулалт} \\
\hline
Jest & Backend unit болон integration туршилт \\
\hline
Supertest & HTTP endpoint туршилт \\
\hline
Vitest & Frontend unit туршилт \\
\hline
React Testing Library & Component туршилт \\
\hline
Playwright & E2E туршилт \\
\hline
Postman / Newman & API туршилт, CI pipeline \\
\hline
k6 & Load ба performance туршилт \\
\hline
\end{tabular}
\end{table}

\subsection{Гол туршилтын кейсүүд}

\begin{table}[h!]
\centering
\caption{Баталгаажуулалтын туршилтын кейсүүд}
\label{tab:test_auth}
\begin{tabular}{|c|p{4cm}|p{3cm}|c|}
\hline
\textbf{TC-ID} & \textbf{Нэр} & \textbf{Оролт} & \textbf{Хүлээгдсэн үр дүн} \\
\hline
TC-001 & Амжилттай бүртгэл & Зөв email, нууц үг & 201, JWT token \\
\hline
TC-002 & Давхардсан email & Бүртгэлтэй email & 409 Conflict \\
\hline
TC-003 & Буруу нууц үг & Буруу password & 401 Unauthorized \\
\hline
TC-004 & Token refresh & Хүчинтэй refresh token & 200, шинэ token \\
\hline
TC-005 & Дууссан token & Expired access token & 401 Unauthorized \\
\hline
\end{tabular}
\end{table}

\begin{table}[h!]
\centering
\caption{Видео үүсгэлтийн туршилтын кейсүүд}
\label{tab:test_video}
\begin{tabular}{|c|p{4cm}|p{3cm}|c|}
\hline
\textbf{TC-ID} & \textbf{Нэр} & \textbf{Оролт} & \textbf{Хүлээгдсэн үр дүн} \\
\hline
TC-010 & Амжилттай видео үүсгэлт & Зөв сэдэв, JWT & 200, jobId \\
\hline
TC-011 & Хязгаар хэтэрсэн & Free хэрэглэгч, хязгаар дүүрсэн & 429 Too Many Requests \\
\hline
TC-012 & Скрипт алдаа & Groq API timeout & Gemini fallback \\
\hline
TC-013 & SSE явц & Видео үүсгэлт явцад & Event stream \\
\hline
TC-014 & Studio re-render & Scene засварлаж дахин гаргах & Шинэ видео URL \\
\hline
\end{tabular}
\end{table}

\subsection{Code coverage үр дүн}

\begin{table}[h!]
\centering
\caption{Code coverage үр дүн}
\label{tab:coverage}
\begin{tabular}{|l|c|c|c|c|}
\hline
\textbf{Модуль} & \textbf{Statement} & \textbf{Branch} & \textbf{Function} & \textbf{Line} \\
\hline
Auth service & 94\% & 91\% & 100\% & 94\% \\
\hline
Video service & 88\% & 85\% & 92\% & 88\% \\
\hline
Scene service & 91\% & 88\% & 95\% & 91\% \\
\hline
AI service & 82\% & 78\% & 88\% & 82\% \\
\hline
FFmpeg pipeline & 87\% & 84\% & 90\% & 87\% \\
\hline
\textbf{Нийт дундаж} & \textbf{88.4\%} & \textbf{85.2\%} & \textbf{93\%} & \textbf{88.4\%} \\
\hline
\end{tabular}
\end{table}

\subsection{Гүйцэтгэлийн туршилт}

k6 хэрэгслээр дараах үзүүлэлтүүдийг хэмжсэн:

\begin{table}[h!]
\centering
\caption{Гүйцэтгэлийн туршилтын үр дүн}
\label{tab:perf}
\begin{tabular}{|l|c|c|c|}
\hline
\textbf{Үзүүлэлт} & \textbf{Зорилт} & \textbf{Үр дүн} & \textbf{Статус} \\
\hline
API response P50 & < 200ms & 145ms & \checkmark \\
\hline
API response P95 & < 500ms & 380ms & \checkmark \\
\hline
API response P99 & < 1000ms & 720ms & \checkmark \\
\hline
Видео үүсгэх (60 сек) & < 5 минут & 3.8 минут & \checkmark \\
\hline
Нэгэн зэрэг хэрэглэгч & 50+ & 50 & \checkmark \\
\hline
Error rate & < 1\% & 0.3\% & \checkmark \\
\hline
\end{tabular}
\end{table}

\section{Илэрсэн алдааг засварлах}

Туршилтын явцад илэрсэн гол алдаа болон тэдгээрийн шийдлийг дараах хүснэгтэд харуулав:

\begin{table}[h!]
\centering
\caption{Илэрсэн алдаа ба шийдэл}
\label{tab:bugs}
\begin{tabular}{|c|p{4cm}|p{4cm}|c|}
\hline
\textbf{ID} & \textbf{Алдаа} & \textbf{Шийдэл} & \textbf{Ноцтой байдал} \\
\hline
BUG-001 & Groq API rate limit-д хүрэхэд систем зогсдог & Gemini fallback логик нэмсэн & Өндөр \\
\hline
BUG-002 & FFmpeg-д Монгол үсэг subtitle-д газдаг & UTF-8 font нэмж тохируулсан & Өндөр \\
\hline
BUG-003 & SSE холболт таслагдсан тохиолдолд явц алдагдана & Reconnect логик нэмсэн & Дунд \\
\hline
BUG-004 & Зэрэг олон re-render хүсэлтэд scene дараалал буруудна & Request queue хэрэгжүүлсэн & Дунд \\
\hline
BUG-005 & Refresh token хугацаа дуусахад хэрэглэгч тасарна & Auto refresh механизм нэмсэн & Дунд \\
\hline
\end{tabular}
\end{table}

\section{Үр дүнг баримтжуулах}

\subsection{API documentation}

Swagger/OpenAPI 3.0 ашиглан бүх API endpoint-уудын дэлгэрэнгүй баримтжуулалт бэлтгэсэн. Swagger UI-г \code{/api/docs} замаар хандах боломжтой.

\subsection{Code documentation}

TypeDoc ашиглан TypeScript кодын JSDoc коментоос автоматаар HTML баримтжуулалт үүсгэсэн. Бүх \code{Service}, \code{Controller}, \code{Repository} классуудад коментыг нэмсэн.

\subsection{Хэрэглэгчийн гарын авлага}

MkDocs + Material theme ашиглан дараах хэсгүүдтэй хэрэглэгчийн гарын авлага бэлтгэв:

\begin{itemize}
\item Системд бүртгүүлэх, нэвтрэх
\item Видео үүсгэх алхам алхмаар
\item Studio засварлагч хэрхэн ашиглах
\item Projects зохион байгуулах
\item Subscription болон эрхийн тайлбар
\end{itemize}

\subsection{Системийн ажиллагааны үр дүн}

\begin{table}[h!]
\centering
\caption{Системийн ажиллагааны эцсийн үнэлгээ}
\label{tab:final_eval}
\begin{tabular}{|l|c|c|}
\hline
\textbf{Шаардлага} & \textbf{Зорилт} & \textbf{Биелэлт} \\
\hline
Видео үүсгэх хугацаа & < 5 минут & 3.8 минут \\
\hline
API хариулах хурд P95 & < 500ms & 380ms \\
\hline
Code coverage & > 80\% & 88.4\% \\
\hline
Монгол хэл дэмжлэг & Бүрэн & Бүрэн \\
\hline
Studio засварлагч & Бүрэн & Бүрэн \\
\hline
RBAC хэрэгжилт & Бүрэн & Бүрэн \\
\hline
REST API endpoint & 35+ & 40+ \\
\hline
\end{tabular}
\end{table}

\section{Бүлгийн дүгнэлт}

Энэхүү бүлэгт Generative AI ашиглан богино хэмжээний видео контентыг автоматаар үүсгэх, үнэлэх системийн загварчлал ба хөгжүүлэлт болон туршилтын ажлыг бүрэн авч үзэв. Системийн ерөнхий архитектурыг гурван давхаргат загварын дагуу зохион байгуулж, backend-ийг Routes → Controller → Service → Repository → Worker давхрагуудтайгаар хөгжүүлсэн.

Өгөгдлийн сангийн схемийг 8 үндсэн хүснэгт, 7 төрлийн холбоос, 8 индекстэйгээр зохиомжилж 3NF хүртэл нормчилсан. Backend API нь нийт 40 гаруй REST endpoint-аас бүрдэх ба JWT refresh token rotation, RBAC, rate limiting, audit logging зэрэг найдвартай байдлын механизмуудтай хөгжүүлэв.

Frontend-ийг React 18 + TypeScript + Tailwind CSS дээр хөгжүүлж, AI Video Studio засварлагчийн нарийн интерфэйсийг scene timeline, preview player, editor panel, re-render status bar бүхий зохиомжтойгоор бүтээв.

AI үйлчилгээний интеграцид Groq + Llama 3.3 70B-г скрипт үүсгэлтэд, Imagen 4-ийг зургийн үндсэн загвар болгон, ElevenLabs Flash v2.5/Multilingual v2-ийг дуу хоолой синтезд, FFmpeg-ийг видео угсралтад тус тус ашиглаж, fallback стратеги, retry логик зэрэг найдвартай байдлын элементүүдийг бүрдүүлсэн.

Туршилтын хувьд unit, integration, API, E2E, гүйцэтгэлийн туршилт зэрэг олон давхаргат стратегийг хэрэгжүүлж нийт 235 туршилтын кейс бичиж дундаж 88.4\% code coverage хүрсэн. Гүйцэтгэлийн зорилтот үзүүлэлтүүд (60 секундын видеог 5 минутын дотор үүсгэх, P95 < 500ms) бүгд хангагдсан.
